\chapter{Type Safety, Kontrak API, dan Desain Layanan}
\label{ch:kontrak}

\section*{Tujuan Pembelajaran}
\begin{itemize}
  \item Memformalkan satu API (Application Programming Interface) sebagai
        kontrak OpenAPI, lalu meng-\textit{generate} tipe kliennya dari
        kontrak tersebut.
  \item Membedakan validasi \textit{runtime} dari pemeriksaan
        \textit{compile-time}, lalu menunjukkan kesalahan yang hanya
        tertangkap masing-masing.
  \item Memisahkan perubahan kontrak yang aman dari yang merusak klien lama,
        lalu membuktikannya pada server yang berjalan.
\end{itemize}

\section{Pendahuluan}

API biasanya tidak dirancang, melainkan tumbuh. Satu alamat ditambahkan,
satu field disisipkan, dan satu nama diubah karena nama lamanya terasa
kurang tepat. Selama yang memakai API itu hanya penulisnya sendiri, caranya
bekerja. Begitu ada klien lain, tiap perubahan menjadi taruhan.

Bab ini memindahkan API dari kebiasaan menjadi kontrak. Kontraknya ditulis
sebagai dokumen OpenAPI, lalu dokumen itu dipakai tiga kali: untuk
meng-\textit{generate} tipe di sisi klien, untuk memvalidasi isi permintaan
di batas sistem, dan untuk menguji apakah jawaban server masih sesuai
janjinya. Bagian kedua bab membahas keputusan desain yang menyertainya, yaitu
bentuk \textit{error}, \textit{pagination}, idempotensi, dan pembatasan laju.

Demonya satu API katalog berisi 48 buku, dilayani dalam dua versi sekaligus,
dengan satu \textit{gateway} di depannya. Referensi utamanya adalah OpenAPI
3.1 untuk bentuk kontraknya \cite{openapi-31}, JSON Schema 2020-12 untuk
skemanya \cite{jsonschema-2020}, RFC 9457 untuk bentuk \textit{error}
\cite{rfc9457}, RFC 9110 untuk arti metode yang idempoten \cite{rfc9110},
dan disertasi Fielding untuk \textit{constraint} REST (Representational
State Transfer) \cite{fielding2000rest}.

\section{Kontrak sebagai Sumber Kebenaran}

Kontrak adalah satu dokumen yang menyebutkan alamat, parameter, bentuk isi,
dan bentuk jawaban sebuah API. Pada pengembangan \textit{contract-first},
dokumen itu ditulis lebih dahulu, dan kode mengikutinya. Arah sebaliknya juga
mungkin, yaitu kontrak dibuat dari kode, tetapi arah itu membuat kontrak
selalu tertinggal satu langkah.

Analoginya seperti gambar kerja sebuah bangunan. Tukang tidak menebak lebar
pintu dari dinding yang sudah berdiri, dan gambar kerjanya tidak digambar
ulang sesudah pintunya dipasang. Gambar itu yang menjadi acuan, dan perubahan
apa pun dibahas di gambarnya lebih dahulu.

Gambar \ref{fig:kontrak} memperlihatkan satu kontrak yang dipakai empat kali
pada demo bab ini.

\begin{figure}[!htb]
  \centering
  \begin{tikzpicture}[
    font=\scriptsize,
    kotak/.style={draw, rounded corners, minimum height=0.8cm,
                  minimum width=2.7cm, align=center, font=\scriptsize\bfseries},
    panah/.style={-{Stealth[length=2mm]}, thick}
  ]
    \node[kotak, fill=praditagreen!15, minimum height=1.0cm] (kontrak) at (0,0) {kontrak.json\\{\tiny\mdseries OpenAPI 3.1, 5 skema}};
    \node[kotak, fill=orange!12] (tipe) at (5.0,1.6) {Tipe TypeScript\\{\tiny\mdseries 50 baris, di-generate}};
    \node[kotak, fill=blue!8] (validasi) at (5.0,0.5) {Validasi di batas\\{\tiny\mdseries jsonschema}};
    \node[kotak, fill=blue!8] (uji) at (5.0,-0.6) {Uji kontrak\\{\tiny\mdseries jawaban server}};
    \node[kotak, fill=gray!35] (diff) at (5.0,-1.7) {Beda versi\\{\tiny\mdseries aman atau merusak}};

    \draw[panah] (kontrak.east) -- ++(0.4,0) |- (tipe.west);
    \draw[panah] (kontrak.east) -- ++(0.4,0) |- (validasi.west);
    \draw[panah] (kontrak.east) -- ++(0.4,0) |- (uji.west);
    \draw[panah] (kontrak.east) -- ++(0.4,0) |- (diff.west);
  \end{tikzpicture}
  \caption{Satu kontrak dipakai empat kali: tipe, validasi, pengujian, dan
           perbandingan versi}
  \label{fig:kontrak}
\end{figure}

Listing \ref{lst:kontrak12} memperlihatkan satu skema di dalam kontrak itu.
Dua hal pada skema tersebut menentukan perilaku seluruh demo:
\texttt{required} menyebut field yang wajib ada, dan
\texttt{additionalProperties} bernilai \texttt{false} menolak field yang
tidak disebutkan.

\begin{lstlisting}[language=json, caption={Satu skema di dalam kontrak, tersedia di \berkas{sources/chapter12/kontrak.json}}, label={lst:kontrak12}]
      "PesananBaru": {
        "type": "object",
        "required": ["id_buku", "jumlah"],
        "additionalProperties": false,
        "properties": {
          "id_buku": { "type": "integer", "minimum": 1 },
          "jumlah": { "type": "integer", "minimum": 1, "maximum": 10 }
        }
      }
\end{lstlisting}

\section{Validasi di Batas Sistem}

Validasi di batas sistem berarti satu pemeriksaan di pintu masuk, bukan
pemeriksaan yang diulang di tiap fungsi. Isi permintaan diperiksa terhadap
skema di kontrak, dan sesudah lolos, kode di belakangnya boleh
memperlakukannya sebagai data yang bentuknya sudah pasti.

Listing \ref{lst:validasi12} memperlihatkan seluruh validasinya. Panjangnya
delapan baris, karena aturannya tidak ditulis di sana, melainkan dibaca dari
kontrak.

\begin{lstlisting}[style=PythonStyle, caption={Validasi di batas sistem memakai skema dari kontrak, tersedia di \berkas{sources/chapter12/aplikasi.py}}, label={lst:validasi12}]
def periksa_isi(versi, isi):
  """Memvalidasi isi permintaan terhadap skema PesananBaru di kontrak.

  Inilah validasi di batas sistem: satu pemeriksaan di pintu masuk, memakai
  skema yang sama dengan yang dipakai meng-generate tipe kliennya.
  """
  try:
    jsonschema.validate(isi, skema(versi, "PesananBaru"))
  except jsonschema.ValidationError as kesalahan:
    return kesalahan.message
  return None
\end{lstlisting}

Tabel \ref{tab:validasi} memuat hasil pengiriman lima isi permintaan ke
alamat \texttt{/v1/pesanan}, termasuk yang sengaja salah. Statusnya
dibandingkan dengan yang dijanjikan kontrak.

\begin{table}[!htb]
  \centering
  \caption{Lima isi permintaan di batas sistem, hasil pengukuran nyata}
  \label{tab:validasi}
  \begin{tabularx}{\textwidth}{Xll}
    \hline
    \textbf{Isi yang dikirim} & \textbf{Status} & \textbf{Alasan dari server} \\
    \hline
    sesuai kontrak & 201 & diterima \\
    \hline
    field wajib \texttt{id\_buku} hilang & 422 & isi tidak sesuai \\
    \hline
    \texttt{id\_buku} bertipe string & 422 & isi tidak sesuai \\
    \hline
    \texttt{jumlah} bernilai 99 & 422 & isi tidak sesuai \\
    \hline
    ada field asing \texttt{catatan} & 422 & isi tidak sesuai \\
    \hline
  \end{tabularx}
\end{table}

Baris terakhir sering mengejutkan. Field asing ditolak, bukan diabaikan,
karena \texttt{additionalProperties} bernilai \texttt{false}. Pilihan itu
punya dua akibat. Pertama, salah tulis nama field ketahuan seketika, bukan
berujung pada data yang hilang tanpa pesan. Kedua, menambah field baru
menjadi perubahan kontrak yang harus disengaja.

\section{Tipe Hasil Generate dan Validasi Compile-Time}

Validasi \textit{runtime} memeriksa data yang benar-benar tiba, dan baru
berjalan ketika programnya jalan. Pemeriksaan \textit{compile-time}
memeriksa kode sebelum dijalankan, dan tidak pernah melihat data sungguhan.
Keduanya diperlukan, dan Tabel \ref{tab:dua-validasi} memisahkan tugasnya.

\begin{table}[!htb]
  \centering
  \caption{Validasi \textit{runtime} lawan pemeriksaan \textit{compile-time}}
  \label{tab:dua-validasi}
  \begin{tabularx}{\textwidth}{lXX}
    \hline
    \textbf{Hal} & \textbf{Runtime} & \textbf{Compile-time} \\
    \hline
    Yang diperiksa & data yang tiba & kode yang ditulis \\
    \hline
    Kapan & saat permintaan dilayani & sebelum program jalan \\
    \hline
    Menangkap & isi permintaan yang salah & pemakaian field yang tidak ada \\
    \hline
    Tidak menangkap & salah tulis di kode klien & data yang dikirim klien lain \\
    \hline
    Pada demo ini & \texttt{jsonschema} di server & \texttt{tsc --strict} di klien \\
    \hline
  \end{tabularx}
\end{table}

Tipe di sisi klien tidak ditulis tangan. Generator membacanya dari kontrak,
lalu menuliskannya sebagai tipe TypeScript, termasuk field opsional dan
daftar nilai yang diperbolehkan. Dari kontrak berisi lima skema dan dua
jalur, hasilnya 50 baris tipe. Listing \ref{lst:tipe12} memuat sebagiannya.

\begin{lstlisting}[language=TypeScript, caption={Tipe hasil generate dari kontrak, tersedia di \berkas{sources/chapter12/tipe-katalog.ts}}, label={lst:tipe12}]
export interface HalamanBuku {
  item: Buku[];
  halaman: number;
  per_halaman: number;
  total: number;
  berikutnya?: string;
}

export interface PesananBaru {
  id_buku: number;
  jumlah: number;
}

export interface Pesanan {
  id: string;
  id_buku: number;
  jumlah: number;
  status: "diterima" | "diproses";
}
\end{lstlisting}

Dua hal pada potongan itu datang langsung dari kontrak. Tanda tanya pada
\texttt{berikutnya} datang dari daftar \texttt{required}, sehingga klien wajib
memeriksanya sebelum dipakai. Daftar nilai pada \texttt{status} datang dari
\texttt{enum}, sehingga membandingkan status dengan nilai yang tidak ada di
kontrak pun tertolak compiler.

Gunanya terlihat ketika kontraknya berubah. Pada versi 2, field
\texttt{harga} berganti nama menjadi \texttt{harga\_rupiah}. Klien lama yang
masih membaca \texttt{harga} langsung ditolak, seperti pada Listing
\ref{lst:tsc12}.

\begin{lstlisting}[language=bash, caption={Kesalahan yang tertangkap sebelum program jalan, memakai \berkas{sources/chapter12/pakai-tipe-salah.ts}}, label={lst:tsc12}]
../chapter06/node_modules/.bin/tsc --noEmit --strict --ignoreConfig \
  pakai-tipe-salah.ts

# Keluarannya:
# pakai-tipe-salah.ts(15,15): error TS2339: Property 'harga' does not
# exist on type 'Buku'.
\end{lstlisting}

\section{Model Error, Pagination, dan Filtering}

Bentuk \textit{error} yang seragam membuat klien cukup menulis satu
penanganan kesalahan. Demo bab ini memakai Problem Details, yaitu
\texttt{application/problem+json} dengan field \texttt{type},
\texttt{title}, \texttt{status}, dan \texttt{detail} \cite{rfc9457}. Nilai
\texttt{detail} berisi pesan dari validator, sehingga klien tahu field mana
yang salah tanpa menebak.

\textit{Pagination} dan \textit{filtering} ditulis di kontrak sebagai
parameter, lengkap dengan batasnya. Parameter \texttt{per\_halaman} punya
nilai bawaan 10 dan batas 50, dan batas itu bukan hiasan: permintaan
\texttt{per\_halaman=99} dijawab 422 beserta alasannya. Tabel
\ref{tab:paging} memuat ketiga jawaban yang diukur.

\begin{table}[!htb]
  \centering
  \caption{\textit{Pagination}, \textit{filtering}, dan batasnya, hasil
           pengukuran nyata pada katalog 48 buku}
  \label{tab:paging}
  \begin{tabularx}{\textwidth}{Xll}
    \hline
    \textbf{Permintaan} & \textbf{Status} & \textbf{Isi jawaban} \\
    \hline
    \texttt{/v1/buku} & 200 & 10 item, total 48, ada \texttt{berikutnya} \\
    \hline
    \texttt{/v1/buku?halaman=2} & 200 & 10 item, halaman 2 \\
    \hline
    \texttt{/v1/buku?per\_halaman=99} & 422 & Problem Details, batas 50 \\
    \hline
  \end{tabularx}
\end{table}

Ketiga jawaban itu divalidasi terhadap skema di kontrak, dan ketiganya
sesuai. Pengujian seperti itu murah dan sering menyelamatkan: jawaban yang
bentuknya menyimpang dari kontrak ketahuan di mesin pengembang, bukan di
aplikasi klien.

\section{Idempotency}

Metode \texttt{GET} tidak mengubah apa pun, sehingga mengulanginya aman.
\texttt{POST} tidak punya jaminan itu \cite{rfc9110}. Pada jaringan yang
putus sesudah permintaan terkirim tetapi sebelum jawabannya tiba, klien tidak
tahu pesanannya sudah masuk atau belum, dan satu-satunya pilihannya adalah
mengirim ulang.

Analoginya seperti menyetor uang kuliah di loket. Mahasiswa membawa nomor
setoran, dan petugas mencatat nomor itu. Ketika mahasiswa yang sama datang
lagi dengan nomor yang sama karena struknya hilang, petugas tidak mencatat
setoran kedua, melainkan mencetak ulang struk yang sudah ada.

Nomor itulah \textit{idempotency key}. Listing \ref{lst:idempoten12}
memperlihatkan penanganannya, termasuk kasus yang mudah terlewat: kunci yang
sama dipakai untuk isi yang berbeda.

\begin{lstlisting}[style=PythonStyle, caption={Satu kunci, satu pesanan, tersedia di \berkas{sources/chapter12/aplikasi.py}}, label={lst:idempoten12}]
def buat_pesanan(versi, isi, kunci):
  """Membuat satu pesanan, atau mengembalikan pesanan lama untuk kunci sama.

  Mengembalikan pasangan (pesanan, status). Status 200 berarti kunci yang
  sama sudah pernah dipakai dengan isi yang sama, sehingga tidak ada pesanan
  baru yang dibuat.
  """
  sidik = json.dumps(isi, sort_keys=True)
  if kunci in IDEMPOTENSI:
    sidik_lama, id_pesanan = IDEMPOTENSI[kunci]
    if sidik_lama != sidik:
      return None, 409
    return PESANAN[id_pesanan], 200
  id_pesanan = f"psn-{len(PESANAN) + 1:04d}"
  pesanan = {"id": id_pesanan, "id_buku": isi["id_buku"],
             "jumlah": isi["jumlah"], "status": "diterima"}
  PESANAN[id_pesanan] = pesanan
  IDEMPOTENSI[kunci] = (sidik, id_pesanan)
  return pesanan, 201
\end{lstlisting}

Tabel \ref{tab:idempoten} memuat hasilnya pada server yang baru menyala.

\begin{table}[!htb]
  \centering
  \caption{Tiga permintaan dengan kunci idempotensi, hasil pengukuran nyata}
  \label{tab:idempoten}
  \begin{tabularx}{\textwidth}{Xll}
    \hline
    \textbf{Permintaan} & \textbf{Status} & \textbf{Hasil} \\
    \hline
    kunci baru & 201 & pesanan \texttt{psn-0002} dibuat \\
    \hline
    kunci sama, isi sama & 200 & \texttt{psn-0002} dikembalikan, bukan pesanan baru \\
    \hline
    kunci sama, isi beda & 409 & ditolak, kunci sudah dipakai \\
    \hline
  \end{tabularx}
\end{table}

Perhatikan status pada baris kedua. Jawabannya 200, bukan 201, karena tidak
ada yang dibuat. Selisih satu angka itu yang membedakan "pesanan Anda sudah
masuk" dari "pesanan Anda masuk dua kali".

\section{Kompatibilitas dan Versioning}

Perubahan kontrak dibagi dua. Perubahan yang aman hanya menambah hal
opsional, sehingga klien lama tidak merasakan apa-apa. Perubahan yang merusak
menghapus atau mengganti hal yang wajib, dan klien lama berhenti bekerja
walaupun servernya sehat.

Pembagian itu dapat diperiksa mesin, bukan diingat orang. Demo bab ini
membandingkan kedua kontrak lalu menamai tiap perubahannya, seperti pada
Tabel \ref{tab:beda-versi}.

\begin{table}[!htb]
  \centering
  \caption{Perubahan kontrak versi 1 ke versi 2, hasil perbandingan otomatis}
  \label{tab:beda-versi}
  \begin{tabularx}{\textwidth}{Xl}
    \hline
    \textbf{Perubahan} & \textbf{Akibat} \\
    \hline
    \texttt{Buku.harga} wajib dihapus & merusak klien lama \\
    \hline
    \texttt{Buku.harga\_rupiah} wajib ditambah & merusak klien lama \\
    \hline
    \texttt{Buku.penerbit} opsional ditambah & aman \\
    \hline
    \texttt{Pesanan.status} enum menerima \texttt{dibatalkan} & merusak klien lama \\
    \hline
    \texttt{PesananBaru.jumlah} batas naik ke 20 & aman bagi klien lama \\
    \hline
  \end{tabularx}
\end{table}

Baris terakhir perlu dibaca dua arah. Batas yang dilonggarkan aman bagi klien
lama, karena permintaan lamanya tetap sah. Arah sebaliknya tidak aman: klien
versi 2 yang memesan 15 buku ditolak 422 oleh versi 1, sebab di sana batasnya
masih 10.

Tabel \ref{tab:uji-versi} memuat pengujiannya pada server yang berjalan.
Keempat baris itu diukur sekali jalan, dengan skema versi 1 sebagai pemeriksa.

\begin{table}[!htb]
  \centering
  \caption{Klien lama dan klien baru diuji ke kedua versi, hasil pengukuran
           nyata}
  \label{tab:uji-versi}
  \begin{tabularx}{\textwidth}{Xll}
    \hline
    \textbf{Yang diuji} & \textbf{Status} & \textbf{Hasil} \\
    \hline
    klien v1 membaca \texttt{/v1/buku} & 200 & jalan \\
    \hline
    klien v1 membaca \texttt{/v2/buku} & 200 & gagal, \texttt{harga} tidak ada \\
    \hline
    klien v2 memesan 15 ke \texttt{/v2} & 201 & diterima \\
    \hline
    klien v2 memesan 15 ke \texttt{/v1} & 422 & isi tidak sesuai \\
    \hline
  \end{tabularx}
\end{table}

Baris kedua memuat pelajaran terpentingnya. Statusnya 200, dan servernya
tidak salah sedikit pun. Yang rusak adalah kesepakatannya, dan kerusakan itu
hanya terlihat bila ada yang memvalidasi jawaban terhadap kontrak lama. Karena
itu versi lama tetap dilayani di alamatnya sendiri, dan klien dipindahkan
ketika siap, bukan ketika servernya di-\textit{deploy}.

\section{Rate Limiting dan API Gateway}

\textit{Gateway} adalah satu pintu masuk di depan beberapa layanan. Tugasnya
hal-hal yang berlaku bagi semua permintaan: membatasi laju, menempelkan
penanda permintaan, memeriksa penanda identitas, lalu meneruskan sisanya.
Dengan begitu aturan itu ditulis sekali, bukan di tiap layanan.

Pembatasan lajunya memakai \textit{token bucket}. Tiap klien punya ember
berisi token, satu permintaan memakai satu token, dan embernya terisi ulang
dengan laju tetap. Analoginya seperti kupon fotokopi gratis di sekretariat:
jatahnya diisi berkala, boleh dipakai sekaligus atau sedikit-sedikit, dan
begitu habis harus menunggu pengisian berikutnya.

Gambar \ref{fig:ember} memperlihatkan jalannya satu permintaan pada demo ini,
yang embernya berisi 20 token dan terisi 5 token per detik.

\begin{figure}[!htb]
  \centering
  \begin{tikzpicture}[
    font=\scriptsize,
    kotak/.style={draw, rounded corners, minimum height=0.85cm,
                  minimum width=2.3cm, align=center, font=\scriptsize\bfseries},
    panah/.style={-{Stealth[length=2mm]}, thick}
  ]
    \node[kotak, fill=blue!8] (klien) at (0,0) {Klien};
    \node[kotak, fill=orange!12, minimum height=1.1cm] (ember) at (3.3,0) {Gateway\\{\tiny\mdseries ember 20 token}\\{\tiny\mdseries isi 5 per detik}};
    \node[kotak, fill=praditagreen!15] (aplikasi) at (7.0,0.7) {Aplikasi 8090\\{\tiny\mdseries v1 dan v2}};
    \node[kotak, fill=red!22] (tolak) at (7.0,-0.8) {429\\{\tiny\mdseries Retry-After}};

    \draw[panah] (klien) -- (ember);
    \draw[panah] (ember.east) -- ++(0.3,0) |- node[font=\tiny, above, pos=0.75] {token ada} (aplikasi.west);
    \draw[panah] (ember.east) -- ++(0.3,0) |- node[font=\tiny, below, pos=0.75] {token habis} (tolak.west);
  \end{tikzpicture}
  \caption{\textit{Token bucket} di \textit{gateway}: lonjakan pendek
           dilayani, laju rata-rata dibatasi}
  \label{fig:ember}
\end{figure}

Tabel \ref{tab:gateway} memuat hasil pengukurannya. Tiga tahap diukur
berurutan, dan angkanya cocok dengan pengaturan embernya.

\begin{table}[!htb]
  \centering
  \caption{Lonjakan permintaan di \textit{gateway}, hasil pengukuran nyata}
  \label{tab:gateway}
  \begin{tabularx}{\textwidth}{Xllll}
    \hline
    \textbf{Tahap} & \textbf{200} & \textbf{429} & \textbf{Retry-After} & \textbf{Lama} \\
    \hline
    30 permintaan sekaligus & 20 & 10 & 0,19 s & 0,02 s \\
    \hline
    setelah menunggu 2 detik & 10 & 20 & 0,16 s & 0,02 s \\
    \hline
    klien lain, ember sendiri & 20 & 10 & 0,19 s & 0,02 s \\
    \hline
  \end{tabularx}
\end{table}

Ketiga baris itu dapat dihitung sebelum diukur, dan itulah sebabnya layak
dibaca teliti. Baris pertama melayani 20 permintaan karena embernya berisi 20
token. Baris kedua melayani 10, karena menunggu 2 detik mengisi $2 \times 5$
token. Baris ketiga melayani 20 lagi karena embernya milik klien lain, bukan
milik sistem.

Pola BFF (Backend for Frontend) duduk satu lapis di atasnya, dan sudah
dibahas di Bab 8. Bedanya begini. \textit{Gateway} tidak tahu kebutuhan
layarnya dan memperlakukan semua klien sama, sedangkan BFF memang dibuat
untuk satu jenis klien dan boleh menggabungkan beberapa jawaban menjadi satu
\cite{newman-bff}.

\section{Gaya API: REST, GraphQL, dan RPC}

Tiga gaya API dipakai luas, dan ketiganya menjawab pertanyaan yang sama
dengan cara berbeda. REST menyusun API sebagai kumpulan sumber daya beralamat
sendiri, dengan metode HTTP sebagai kata kerjanya \cite{fielding2000rest}.
GraphQL menyediakan satu alamat, dan kliennya yang menyebutkan field apa saja
yang diinginkan \cite{graphql-spec}. RPC (Remote Procedure Call) memandang
API sebagai pemanggilan fungsi, dan gRPC menuliskan kontraknya sebagai berkas
\textit{protocol buffer} \cite{grpc-docs}.

Tabel \ref{tab:gaya-api} membandingkan ketiganya pada hal-hal yang biasanya
menentukan pilihan.

\begin{table}[!htb]
  \centering
  \caption{Tiga gaya API dan \textit{tradeoff}-nya}
  \label{tab:gaya-api}
  \begin{tabularx}{\textwidth}{lXXX}
    \hline
    \textbf{Hal} & \textbf{REST} & \textbf{GraphQL} & \textbf{RPC} \\
    \hline
    Bentuk & sumber daya beralamat & satu alamat, \textit{query} klien & daftar fungsi \\
    \hline
    Kontrak & OpenAPI & skema GraphQL & \textit{protocol buffer} \\
    \hline
    Cache HTTP & langsung dapat dipakai & perlu usaha tambahan & umumnya tidak \\
    \hline
    Ambil data berlebih & sering & jarang & tergantung fungsinya \\
    \hline
    Cocok untuk & API publik & layar dengan kebutuhan beragam & layanan ke layanan \\
    \hline
  \end{tabularx}
\end{table}

Pilihan gayanya jarang menjadi penentu mutu API. Yang menentukan adalah hal
yang sama pada ketiganya, yaitu ada kontrak yang dapat dibaca mesin, bentuk
\textit{error} yang seragam, batas laju yang jelas, dan perubahan versi yang
tidak diam-diam.

\section{Ringkasan}

Kontrak memindahkan kesepakatan API dari ingatan ke dokumen. Pada demo bab
ini, satu dokumen OpenAPI berisi lima skema dan dua jalur dipakai empat kali:
meng-\textit{generate} 50 baris tipe TypeScript, memvalidasi isi permintaan
di batas sistem, memvalidasi jawaban server, dan membandingkan dua versi
kontrak. Tidak ada aturan yang ditulis dua kali, sehingga tidak ada aturan
yang dapat menyimpang sendiri.

Validasi \textit{runtime} dan pemeriksaan \textit{compile-time} menangkap
kesalahan yang berbeda. Dari lima isi permintaan yang diuji, empat ditolak
422, termasuk satu yang hanya menambah field asing. Pemeriksaan
\textit{compile-time} menangkap hal yang tidak pernah dilihat server, yaitu
klien yang membaca field \texttt{harga} pada kontrak versi 2, dan menolaknya
sebagai TS2339 sebelum program jalan.

Idempotensi mengubah pengiriman ulang dari risiko menjadi hal biasa. Kunci
yang sama dengan isi yang sama dijawab 200 beserta pesanan yang sudah ada,
bukan 201 beserta pesanan kedua. Kunci yang sama dengan isi berbeda dijawab
409, karena menerima permintaan itu diam-diam berarti mencatat sesuatu yang
tidak diminta siapa pun.

Kompatibilitas dapat diperiksa mesin. Perbandingan kedua kontrak menamai lima
perubahan: dua menghapus dan menambah field wajib, satu menambah field
opsional yang aman, satu menambah nilai enum, dan satu melonggarkan batas.
Pada server yang berjalan, klien versi 1 membaca jawaban versi 2 dengan
status 200 tetapi gagal memakainya, dan kegagalan seperti itu hanya terlihat
bila ada yang memvalidasinya terhadap kontrak lama.

Pesan utama bab ini adalah tempat aturan disimpan. Aturan yang hidup di
kontrak dapat di-\textit{generate}, divalidasi, dibandingkan, dan diuji.
Aturan yang hidup di kepala penulisnya hanya dapat diingat, dan ingatan tidak
pernah menjadi bagian dari \textit{pipeline}.

\section{Latihan}

Seluruh file yang dipakai pada latihan ini tersedia di
\berkas{sources/chapter12/}, dihitung relatif terhadap akar repositori.
Latihan ini membutuhkan Python 3 beserta \texttt{jsonschema}, dan
\texttt{tsc} yang sudah terpasang di Bab 6. Tidak ada Docker dan tidak ada
pemasangan npm baru.

Seluruh permintaan diarahkan ke aplikasi dan \textit{gateway} yang berjalan
di komputer sendiri, dan sifatnya mengamati, bukan membebani. Status dan
klasifikasi perubahan akan sama pada tiap mahasiswa, sedangkan lama
pengukuran dan nilai \texttt{Retry-After} akan berbeda sedikit.

Penyiapannya dua perintah, seperti pada Listing \ref{lst:c12-siapkan}.

\begin{lstlisting}[language=bash, caption={Menyalakan aplikasi dan gateway, tersedia di \berkas{sources/chapter12/}}, label={lst:c12-siapkan}]
cd sources/chapter12
python3 aplikasi.py &
python3 gateway.py &

# Keluarannya:
# Melayani di http://127.0.0.1:8090/v1/buku dengan 48 buku
# Gateway di http://127.0.0.1:8091/ dengan ember 20 token dan 5 token per detik
\end{lstlisting}

\subsection*{Latihan 1: Meng-generate Tipe dari Kontrak}

Latihan ini membuktikan bahwa tipe klien dapat diturunkan dari kontrak, dan
bahwa perubahan kontrak tertangkap compiler sebelum program jalan.

\begin{enumerate}
  \item Jalankan generatornya untuk kedua kontrak, seperti pada Listing
        \ref{lst:c12-lat1-jalan}, lalu catat jumlah skema dan jumlah baris
        hasilnya.
  \item Periksa seluruh tipe dengan \texttt{tsc} memakai
        \berkas{tsconfig.json}, lalu catat hasilnya.
  \item Periksa \berkas{pakai-tipe-salah.ts} sendiri, lalu catat kode
        kesalahan beserta nomor barisnya.
  \item Tambahkan satu field opsional pada skema \texttt{Buku} di
        \berkas{kontrak.json}, jalankan ulang generatornya, lalu catat
        perubahan pada file tipenya.
  \item Isi Tabel \ref{tab:c12-lat1-hasil}. Tabel \ref{tab:c12-lat1-contoh}
        memperlihatkan contoh yang sudah terisi.
  \item Jawab dua pertanyaan berikut. Pertama, mengapa field yang tidak ada
        di \texttt{required} di-\textit{generate} dengan tanda tanya? Kedua,
        kesalahan apa yang tetap lolos dari \texttt{tsc} tetapi tertangkap
        validasi di server?
\end{enumerate}

\begin{lstlisting}[language=bash, caption={Meng-generate tipe lalu memeriksanya, tersedia di \berkas{sources/chapter12/buat\_tipe.py}}, label={lst:c12-lat1-jalan}]
python3 buat_tipe.py kontrak.json tipe-katalog.ts
python3 buat_tipe.py kontrak-v2.json tipe-katalog-v2.ts
../chapter06/node_modules/.bin/tsc --noEmit -p tsconfig.json

# Keluarannya:
# 5 skema dan 2 jalur di-generate menjadi tipe-katalog.ts
# 5 skema dan 2 jalur di-generate menjadi tipe-katalog-v2.ts
\end{lstlisting}

\begin{table}[!htb]
  \centering
  \caption{Contoh hasil Latihan 1 yang sudah terisi}
  \label{tab:c12-lat1-contoh}
  \begin{tabularx}{\textwidth}{Xll}
    \hline
    \textbf{Yang diperiksa} & \textbf{Hasil} & \textbf{Keterangan} \\
    \hline
    Jumlah skema dan jalur & 5 dan 2 & dari \berkas{kontrak.json} \\
    \hline
    Baris tipe hasil generate & 50 & versi 2 menjadi 51 \\
    \hline
    \texttt{tsc} pada seluruh tipe & lolos & tanpa kesalahan \\
    \hline
    \texttt{tsc} pada file salah & TS2339 & baris 15, field \texttt{harga} \\
    \hline
  \end{tabularx}
\end{table}

\begin{table}[!htb]
  \centering
  \caption{Lembar isian hasil pengamatan Latihan 1}
  \label{tab:c12-lat1-hasil}
  \begin{tabularx}{\textwidth}{Xll}
    \hline
    \textbf{Yang diperiksa} & \textbf{Hasil} & \textbf{Keterangan} \\
    \hline
    Jumlah skema dan jalur & \dots & \dots \\
    \hline
    Baris tipe hasil generate & \dots & \dots \\
    \hline
    \texttt{tsc} pada seluruh tipe & \dots & \dots \\
    \hline
    \texttt{tsc} pada file salah & \dots & \dots \\
    \hline
  \end{tabularx}
\end{table}

\subsection*{Latihan 2: Validasi di Batas dan Idempotensi}

Latihan ini membuktikan bahwa satu pemeriksaan di pintu masuk cukup untuk
menolak lima bentuk isi yang salah, dan bahwa pengiriman ulang tidak membuat
pesanan kedua.

\begin{enumerate}
  \item Nyalakan ulang aplikasinya, agar nomor pesanannya dimulai dari awal.
  \item Jalankan \berkas{uji-kontrak.py}, seperti pada Listing
        \ref{lst:c12-lat2-jalan}, lalu catat sebelas barisnya.
  \item Kirim satu permintaan dengan kunci idempotensi pilihan sendiri dua
        kali, lalu catat kedua statusnya dan nomor pesanannya.
  \item Kirim permintaan ketiga dengan kunci yang sama tetapi
        \texttt{id\_buku} berbeda, lalu catat statusnya.
  \item Isi Tabel \ref{tab:c12-lat2-hasil}. Tabel \ref{tab:c12-lat2-contoh}
        memperlihatkan contoh yang sudah terisi.
  \item Jawab tiga pertanyaan berikut. Pertama, mengapa pengiriman kedua
        dijawab 200, bukan 201? Kedua, mengapa field asing ditolak, bukan
        diabaikan? Ketiga, apa yang terjadi bila kunci idempotensinya dibuat
        dari isi permintaan, bukan dibuat klien?
\end{enumerate}

\begin{lstlisting}[language=bash, caption={Menguji validasi dan idempotensi, tersedia di \berkas{sources/chapter12/uji-kontrak.py}}, label={lst:c12-lat2-jalan}]
python3 uji-kontrak.py

# Keluarannya:
#   Yang diuji                   Status  Hasil
#   GET /v1/buku                    200  sesuai skema
#   sesuai kontrak                  201  sesuai, diterima
#   field asing                     422  sesuai, isi tidak sesuai
#   kunci baru                      201  pesanan psn-0002 dibuat
#   kunci sama, isi sama            200  pesanan psn-0002 dikembalikan, sama
#   kunci sama, isi beda            409  kunci idempotensi dipakai ulang
\end{lstlisting}

\begin{table}[!htb]
  \centering
  \caption{Contoh hasil Latihan 2 yang sudah terisi}
  \label{tab:c12-lat2-contoh}
  \begin{tabularx}{\textwidth}{Xll}
    \hline
    \textbf{Yang dikirim} & \textbf{Status} & \textbf{Hasil} \\
    \hline
    isi sesuai kontrak & 201 & diterima \\
    \hline
    field wajib hilang & 422 & isi tidak sesuai \\
    \hline
    field asing \texttt{catatan} & 422 & isi tidak sesuai \\
    \hline
    kunci sama, isi sama & 200 & pesanan yang sama dikembalikan \\
    \hline
    kunci sama, isi beda & 409 & ditolak \\
    \hline
  \end{tabularx}
\end{table}

\begin{table}[!htb]
  \centering
  \caption{Lembar isian hasil pengamatan Latihan 2}
  \label{tab:c12-lat2-hasil}
  \begin{tabularx}{\textwidth}{Xll}
    \hline
    \textbf{Yang dikirim} & \textbf{Status} & \textbf{Hasil} \\
    \hline
    isi sesuai kontrak & \dots & \dots \\
    \hline
    field wajib hilang & \dots & \dots \\
    \hline
    field asing \texttt{catatan} & \dots & \dots \\
    \hline
    kunci sama, isi sama & \dots & \dots \\
    \hline
    kunci sama, isi beda & \dots & \dots \\
    \hline
  \end{tabularx}
\end{table}

\subsection*{Latihan 3: Memisahkan Perubahan Aman dari yang Merusak}

Latihan ini membuktikan bahwa jenis perubahan kontrak dapat diputuskan mesin,
dan bahwa klien lama dapat rusak walaupun server menjawab 200.

\begin{enumerate}
  \item Jalankan \berkas{uji-kompatibilitas.py}, seperti pada Listing
        \ref{lst:c12-lat3-jalan}, lalu catat lima perubahan beserta
        akibatnya.
  \item Catat keempat baris pengujian pada server, terutama statusnya.
  \item Tambahkan satu field opsional pada skema \texttt{Buku} di
        \berkas{kontrak-v2.json}, jalankan ulang, lalu catat perubahannya di
        tabel perbandingan.
  \item Ganti nama satu field wajib di \berkas{kontrak-v2.json}, jalankan
        ulang, lalu catat berapa baris "merusak klien lama" yang muncul.
  \item Isi Tabel \ref{tab:c12-lat3-hasil}. Tabel \ref{tab:c12-lat3-contoh}
        memperlihatkan contoh yang sudah terisi.
  \item Jawab tiga pertanyaan berikut. Pertama, mengapa mengganti nama satu
        field menghasilkan dua baris perubahan? Kedua, mengapa menambah nilai
        enum merusak klien lama? Ketiga, apa yang harus disiapkan sebelum
        versi 1 dimatikan?
\end{enumerate}

\begin{lstlisting}[language=bash, caption={Membandingkan kontrak lalu mengujinya, tersedia di \berkas{sources/chapter12/uji-kompatibilitas.py}}, label={lst:c12-lat3-jalan}]
python3 uji-kompatibilitas.py

# Keluarannya:
#   Perubahan kontrak v1 ke v2                     Akibat
#   Buku.harga wajib dihapus                       merusak klien lama
#   Buku.penerbit opsional ditambah                aman
#   PesananBaru.jumlah batas naik ke 20            aman bagi klien lama
#
#   Yang diuji pada server             Status  Hasil
#   klien v1 membaca /v2/buku             200  gagal: 'harga' is a required
#   klien v2 memesan 15 ke /v1            422  isi tidak sesuai
\end{lstlisting}

\begin{table}[!htb]
  \centering
  \caption{Contoh hasil Latihan 3 yang sudah terisi}
  \label{tab:c12-lat3-contoh}
  \begin{tabularx}{\textwidth}{Xll}
    \hline
    \textbf{Yang diamati} & \textbf{Hasil} & \textbf{Keterangan} \\
    \hline
    Perubahan terdeteksi & 5 & 2 aman, 3 merusak \\
    \hline
    Klien v1 ke \texttt{/v2/buku} & 200 & gagal dipakai \\
    \hline
    Klien v2 ke \texttt{/v1/pesanan} & 422 & batas lama 10 \\
    \hline
    Ganti nama satu field wajib & 2 baris & dihapus dan ditambah \\
    \hline
  \end{tabularx}
\end{table}

\begin{table}[!htb]
  \centering
  \caption{Lembar isian hasil pengamatan Latihan 3}
  \label{tab:c12-lat3-hasil}
  \begin{tabularx}{\textwidth}{Xll}
    \hline
    \textbf{Yang diamati} & \textbf{Hasil} & \textbf{Keterangan} \\
    \hline
    Perubahan terdeteksi & \dots & \dots \\
    \hline
    Klien v1 ke \texttt{/v2/buku} & \dots & \dots \\
    \hline
    Klien v2 ke \texttt{/v1/pesanan} & \dots & \dots \\
    \hline
    Ganti nama satu field wajib & \dots & \dots \\
    \hline
  \end{tabularx}
\end{table}

\subsection*{Latihan 4: Batas Laju di Gateway}

Latihan ini mengukur perilaku \textit{token bucket}: lonjakan pendek yang
masih dilayani, penolakan sesudah embernya habis, dan pengisian ulangnya.

\begin{enumerate}
  \item Pastikan aplikasi dan \textit{gateway} menyala, lalu jalankan
        \berkas{ukur-gateway.py} seperti pada Listing
        \ref{lst:c12-lat4-jalan}.
  \item Catat jumlah 200 dan 429 pada ketiga tahapnya, beserta
        \texttt{Retry-After} yang diterima.
  \item Hitung sendiri jumlah yang seharusnya dilayani pada tahap kedua dari
        lama menunggu dan laju pengisiannya, lalu bandingkan dengan hasil
        pengukuran.
  \item Ubah \texttt{KAPASITAS\_EMBER} menjadi 5 di \berkas{gateway.py},
        nyalakan ulang \textit{gateway}-nya, lalu ukur sekali lagi.
  \item Isi Tabel \ref{tab:c12-lat4-hasil}. Tabel \ref{tab:c12-lat4-contoh}
        memperlihatkan contoh yang sudah terisi.
  \item Jawab tiga pertanyaan berikut. Pertama, mengapa klien lain tetap
        dilayani 20 permintaan? Kedua, apa akibatnya bila embernya dibuat
        satu untuk seluruh sistem? Ketiga, mengapa penolakannya memakai 429
        beserta \texttt{Retry-After}, bukan 503 tanpa keterangan?
\end{enumerate}

\begin{lstlisting}[language=bash, caption={Mengukur batas laju gateway, tersedia di \berkas{sources/chapter12/ukur-gateway.py}}, label={lst:c12-lat4-jalan}]
python3 ukur-gateway.py

# Keluarannya:
#   Tahap                            200   429  Retry-After  Lama detik
#   30 permintaan sekaligus           20    10         0.19        0.02
#   setelah menunggu 2 detik          10    20         0.16        0.02
#   klien lain, ember sendiri         20    10         0.19        0.02
\end{lstlisting}

\begin{table}[!htb]
  \centering
  \caption{Contoh hasil Latihan 4 yang sudah terisi}
  \label{tab:c12-lat4-contoh}
  \begin{tabularx}{\textwidth}{Xlll}
    \hline
    \textbf{Tahap} & \textbf{200} & \textbf{429} & \textbf{Hitungan sendiri} \\
    \hline
    30 sekaligus & 20 & 10 & kapasitas ember 20 \\
    \hline
    setelah 2 detik & 10 & 20 & $2 \times 5$ token \\
    \hline
    klien lain & 20 & 10 & ember terpisah \\
    \hline
  \end{tabularx}
\end{table}

\begin{table}[!htb]
  \centering
  \caption{Lembar isian hasil pengukuran Latihan 4}
  \label{tab:c12-lat4-hasil}
  \begin{tabularx}{\textwidth}{Xlll}
    \hline
    \textbf{Tahap} & \textbf{200} & \textbf{429} & \textbf{Hitungan sendiri} \\
    \hline
    30 sekaligus & \dots & \dots & \dots \\
    \hline
    setelah 2 detik & \dots & \dots & \dots \\
    \hline
    klien lain & \dots & \dots & \dots \\
    \hline
  \end{tabularx}
\end{table}
